\documentclass[10pt,a4paper]{amsart}
\usepackage[margin=19mm]{geometry}
\usepackage{amsmath,amssymb,mathtools}
\usepackage[scr=rsfs,cal=euler]{mathalfa}
\usepackage{xfrac}
\usepackage[unicode,hidelinks,bookmarks=false]{hyperref}
\newcommand{\ie}{i.e.\ }
\newcommand{\cf}{cf.\ }
\newcommand{\resp}{respectively\ }
\newcommand{\Subsection}{\subsection}
\providecommand{\nobreakdash}{\nobreak\hbox{-}}
\setlength{\emergencystretch}{2em}
\multlinegap0pt
\usepackage{mathtools}
\usepackage{amssymb}
\usepackage[shortlabels]{enumitem}
\usepackage{bm}
\newtheorem{lemma}{Lemma}
\newtheorem{theorem}{Theorem}
\title{On Some Non-Permutations of Quadratic extensions Of Finite Field}
\author{Bidushi Sharma, DHIREN KUMAR BASNET*}
\date{Source: arXiv:2606.14529v2 (CC BY 4.0)}
\begin{document}
\maketitle

\noindent\emph{Source and licence.} On Some Non-Permutations of Quadratic extensions Of Finite Field. Version arXiv:2606.14529v2 (CC BY 4.0), distributed by arXiv under the Creative Commons Attribution 4.0 International licence, http://creativecommons.org/licenses/by/4.0/, as recorded in the arXiv licence metadata for this exact version and shown on the abstract page https://arxiv.org/abs/2606.14529v2. Full licence notice: \texttt{licenses/arxiv-2606.14529-CC-BY-4.0.txt}.\par\smallskip

\noindent\textit{Editorial scope.} Counted unit: the article's theorem t1 (source0.tex lines 753--756). The article's complete original proof of it is delivered with it (source0.tex lines 758--854, one \texttt{proof} environment of 97 lines). No mathematics is written, repaired or re-derived: the statement and the proof are the article's own lines, in the article's own notation, and no retained line is substituted or re-flowed.  Retained uncounted as the source-local context the proof needs: lines 240--251; lines 293--303.  Nothing was omitted from any retained span, so the package carries no omission record.  No retained line contains a citation, so the wrapper carries no bibliography and no bibliography entry.  No retained line contains a figure environment, an image-inclusion command or drawing code, so no image is delivered and the package declares no asset dependency.  Editorial formatting only, in addition: an AMS \texttt{amsart} preamble that restates the article's own declarations and macros used by the retained lines, namely the environments \texttt{lemma}, \texttt{theorem} on separate counters, together with the standard packages those lines need.  The retained environments print in this document as Lemma 1, Theorem 1 in order of appearance, whereas the article prints the same items with the numbering of its own full document.  The complete native source of the article as distributed in the arXiv e-print for this exact version is bundled unchanged as \texttt{sources/202792/source0.tex}.
\begin{lemma}\label{l5}
Let $q=p^k$ and let $c,d \in \mathbb{F}_{q^2}$. Then $M$, the number of zeroes of $
x^q+cx+d$
in $\mathbb{F}_{q^2}$ is given by
$$
M=\begin{cases}
1, & \text{if } c \notin \mu_{q+1},\\[1mm]
q, & \text{if } c \in \mu_{q+1} \text{ and } d=cd^q,\\[1mm]
0, & \text{if } c \in \mu_{q+1} \text{ and } d \neq cd^q.
\end{cases}
$$
\end{lemma}

\begin{lemma}\label{l6}
Let $q$ be odd, $b,d \in \mathbb{F}_{q^2}$ with $b \neq 0$, and $\alpha$ be a primitive element of $\mathbb{F}_{q^2}$. Then the number of zeros of
$x^q+bx^2+d$ in $\mathbb{F}_{q^2}$ is
$$
1-\eta(-b^{q+1}) \sum_{i \in K} (-1)^i,
$$
where
$K=\{\, i \mid 0 \leq i \leq q,\ h(\alpha^{i(q-1)})=0 \,\},
$ and $h(x)=b^q x^3-4b^{q+1}d^q x^2-4b^{q+1}d x+b$.

\end{lemma}

\begin{theorem}\label{t1}
Let $q$ be odd. Then the polynomial
$f(x)=x^q+bx^2+cx+d\in\mathbb{F}_{q^2}[x]$ is not a permutation polynomial of $\mathbb{F}_{q^2}$ if either $b=0$ and $c\in\mu_{q+1}$ or, $b\neq0$.
\end{theorem}

\begin{proof}
First consider the case when $b=0$. Then, by Lemma~\ref{l5}, the result follows directly.

Next assume that $b\neq 0$. Observe that
$
x^q+bx^2+cx+d = g\!\left(x+\dfrac{c}{2b}\right),
$
where
$
g(x)=x^q+bx^2+D,
$
with
$
D=d-\dfrac{c^2}{4b}-\dfrac{c^q}{2b^q}.
$


Since translation preserves the permutation property (and hence the non-permutation property), it suffices to restrict our attention to polynomials of the form $g(x)=x^q+bx^2+D$.
Thus, the equation
$
g(x)+t=0
$
has more than one solution in $\mathbb{F}_{q^2}$ for some $t\in \mathbb{F}_{q^2}$ if and only if
$$
\sum_{i\in K_t} (-1)^i \neq 0,
$$
 where
$
K_t=\left\{\, i \mid 0\leq i\leq q,\; h_t\!\left(\alpha^{\,i(q-1)}\right)=0 \right\},
$
and
$
h_t(x)=b^q x^3-4b^{q+1}(D+t)^q x^2-4b^{q+1}(D+t)x+b.
$
Here, $\alpha$ denotes a primitive element of $\mathbb{F}_{q^2}$. This follows directly from Lemma~\ref{l6}.\\
Notice first that
$
J=\left\{\alpha^{i(q-1)} \mid 0\leq i\leq q \right\}
$
is a subgroup of $\mathbb{F}_{q^2}^{*}$ of order $q+1$, generated by $\alpha^{q-1}$.

We first consider the case $b\in \mathbb{F}_q^{*}$. For
$
t=-D-\dfrac{3}{4b},
$
we have
$
h_t(-1)=0
\quad\text{and}\quad
h_t'(-1)=0.
$
Hence, $-1$ is a root of $h_t(x)=0$ of multiplicity at least $2$. Since $-1\in \mu_{q+1}$ and the product of all roots of $h_t(x)$ is
$
-\dfrac{b}{b^q}=-1,
$
the third root must also be $-1$. Therefore, $-1$ is in fact a root of multiplicity $3$.

Now
$
-1=(\alpha^{q-1})^{\dfrac{q+1}{2}},
$
and therefore
$
K_t=\left\{\dfrac{q+1}{2}\right\}.
$
If $q\equiv 1 \pmod 4$, then $\dfrac{q+1}{2}$ is odd, whereas if $q\equiv 3 \pmod 4$, then $\dfrac{q+1}{2}$ is even. In either case,
$$
\sum_{i\in K_t} (-1)^i \neq 0.
$$
Similarly, assume that $b\in \mathbb{F}_{q^2}\setminus \mathbb{F}_q$ with $\operatorname{Tr}(b)=0$. Take
$
t=\dfrac{2b^{1-q}-1}{4b}-D.
$
Then one readily verifies that $1$ is a root of $h_t(x)=0$ of multiplicity $3$. Consequently,
$
K_t=\{0\}.
$
Therefore,
$$
\sum_{i\in K_t} (-1)^i \neq 0.
$$
Finally, assume that $b\in \mathbb{F}_{q^2}\setminus \mathbb{F}_q$ with $\operatorname{Tr}(b)\neq 0$. For
$
t=\dfrac{1}{4b}-D,
$
the roots of $h_t(x)=0$ are $1$, $-1$, and $\gamma$, where $\gamma\neq ,1-1$. Hence,
$
K_t=\{r,j,k\}.
$

For all possible parity distributions of $r$, $j$ and $k$, we have
$$
\sum_{i\in K_t} (-1)^i \neq 0.
$$
Thus the result follows.

\end{proof}

\end{document}
